\documentclass[10pt,a4paper]{amsart}
\usepackage[margin=19mm]{geometry}
\usepackage{amsmath,amssymb,mathtools}
\usepackage[scr=rsfs,cal=euler]{mathalfa}
\usepackage{xfrac}
\usepackage[unicode,hidelinks,bookmarks=false]{hyperref}
\newcommand{\ie}{i.e.\ }
\newcommand{\cf}{cf.\ }
\newcommand{\resp}{respectively\ }
\newcommand{\Subsection}{\subsection}
\providecommand{\nobreakdash}{\nobreak\hbox{-}}
\setlength{\emergencystretch}{2em}
\multlinegap0pt
\usepackage{bm}
\usepackage{bbm}
\usepackage{dsfont}
\usepackage{xspace}
\usepackage{cancel}
\usepackage{mathtools}
\usepackage{cleveref}
\usepackage{amsthm}
\makeatletter
\@ifundefined{theorem}{\newtheorem{theorem}{Theorem}}{}
\makeatother
\newcommand{\Circ}{\circ}
\newcommand{\distmatrix}{D}
\newcommand{\one}{\vec{1}}
\newcommand{\x}{\vec{x}}
\newcommand{\y}{\vec{y}}
\title[Distance Exceptional Graphs and the Curvature Index]{Distance Exceptional Graphs and the Curvature Index}
\author{Sawyer Jack Robertson, Finn Southerland, Erlang Surya}
\date{Source: arXiv:2511.03719v1 (CC BY 4.0)}
\begin{document}
\maketitle

\noindent\emph{Source and licence.} Distance Exceptional Graphs and the Curvature Index. Version arXiv:2511.03719v1 (CC BY 4.0), distributed under the Creative Commons Attribution 4.0 International licence, as recorded in the arXiv licence metadata for this exact version and shown on the abstract page https://arxiv.org/abs/2511.03719v1. Full rights notice: licenses/arxiv-2511.03719-CC-BY-4.0.txt.\par\smallskip

\noindent\textit{Editorial scope.} Counted unit: the Theorem whose source label is \texttt{thm:merging} in the article (source0.tex lines 1592--1598, source label \texttt{thm:merging}), delivered together with the article's own complete original proof of it (source0.tex lines 1600--1708, one \texttt{proof} environment of 109 lines). No mathematics is written, repaired or re-derived: statement and proof are the article's own text line for line. Required context retained uncounted: none. Every definition, display and label of those lines is retained unchanged. Omitted from this package and not redistributed: 1--1591, 1599--1599, 1709--2536. Nothing inside a retained excerpt is omitted or altered: every retained body is delivered exactly as the article's lines are written, so no omission receipt of any kind accompanies this package. Citations: the retained excerpts carry no citation command at all, so no bibliography entry of the article is transcribed and no reference is redistributed. Reference adaptation: none was needed and none was made; every reference inside the delivered unit resolves inside it, so no unresolved reference or citation is produced. Printed slips kept literal: no printed word, symbol, hypothesis or number of the counted statement or of its proof was altered. Layout and class: the article's own class is \texttt{article} with packages including inputenc, lmodern, fontenc, babel, ifpdf, geometry, xcolor, hyperref, graphicx, amsmath, amsthm, amssymb, enumitem, mathrsfs; the wrapper is the lane's pinned \texttt{amsart} editorial wrapper at 10pt on A4, which restates the theorem-like environments the retained text needs and adds the article's own macro definitions that the retained text needs (\texttt{\textbackslash Circ}, \texttt{\textbackslash distmatrix}, \texttt{\textbackslash one}, \texttt{\textbackslash x}, \texttt{\textbackslash y}), with extra wrapper packages (bm, bbm, dsfont, xspace, cancel, mathtools, cleveref). The delivered numbering is the unit's own. Assets: no retained span contains a figure environment, a graphics-inclusion command, a PDF-page inclusion, a table or a TikZ picture; no image file is copied into this package, the package declares no asset dependency and no image file is redistributed with it. Licence: version arXiv:2511.03719v1 of this article is distributed under the Creative Commons Attribution 4.0 International licence, as recorded in the arXiv licence metadata for this exact version and shown on the abstract page https://arxiv.org/abs/2511.03719v1. A full rights notice is delivered at licenses/arxiv-2511.03719-CC-BY-4.0.txt. Characters: every retained excerpt is pure ASCII, so the delivered engine, XeLaTeX with its default Latin Modern text font, reports no missing character.
    \begin{theorem}[Vertex coalescence]\label{thm:merging}
        Let $G, H$ be connected graphs. Let $u_{0}\in V(G)$ and $v_{0}\in V(H)$ be fixed vertices, and let $G\,\Circ_{u_{0},v_{0}}\,H$ denote the vertex coalescence of $G$ and $H$ at $u_{0}$ and $v_{0}$, respectively. Then
            \begin{align*}
                \iota(G\,\Circ_{u_{0},v_{0}}\,H) &= \iota(G) + \iota(H).
            \end{align*}
        If either $\iota(G) = \infty$ or $\iota(H) = \infty$, then $\iota(G\,\Circ_{u_{0},v_{0}}\,H)=\infty$.
    \end{theorem}

    \begin{proof}
        Fix an enumeration of $V(G)$ and $V(H)$ such that $u_0, v_0$ appear first. Write $G\,\Circ\, H := G\,\Circ_{u_{0},v_{0}}\,H$ for brevity. Assume that $\iota(G),\iota(H)<\infty$. Then there exist vectors $\x_G,\x_H$ with $\one^\top\x_G = \one^\top\x_H = 1$ and $\distmatrix(G)\x_G = \iota(G)\one$, $\distmatrix(H)\x_H = \iota(H)\one$. Write
            \begin{align*}
                \x_G^\top &= \begin{bmatrix}
                    \alpha_G & \y_G^\top
                \end{bmatrix},\quad \x_H^\top = \begin{bmatrix}
                    \alpha_H & \y_H^\top
                \end{bmatrix}, \quad \alpha_G,\alpha_H\in\mathbb{R}.
            \end{align*}
        Order the vertices of $G\,\Circ\,H$ as
            \begin{align*}
                V(G\,\Circ\,H) = \{w\} \cup (V(G)\setminus \{u_{0}\} )\cup (V(H)\setminus \{v_{0}\}),
            \end{align*}
        where $w\in V(G\,\Circ\,H)$ is the vertex appearing after $u_0, v_0$ are identified. Define $s = |V(G)| - 1$, $t = |V(H)| - 1$. Denote by ${\distmatrix(G)}'$ (resp. ${\distmatrix(H)}'$) the principal submatrix of ${\distmatrix(G)}$ (resp. ${\distmatrix(H)}$) indexed by $V(G)\setminus \{u_{0}\}$ (resp. $V(H)\setminus \{v_{0}\}$), obtained by deleting the row and column corresponding to $u_{0}$ (resp. $v_{0}$). Let $\x_{1}\in\mathbb{R}^{s}$ denote the vector of distances between $u_{0}$ and $V(G)\setminus \{u_{0}\}$, and let $\y_{1}\in\mathbb{R}^{t}$ denote the vector of distances between $v_{0}$ and $V(H)\setminus \{v_{0}\}$. With this setup, the distance matrix $\distmatrix = \distmatrix(G\,\Circ\,H)$ takes the form
            \begin{align}\label{eq:d-decomp}
                \distmatrix = \begin{bmatrix}
                    0 & \x_{1}^\top & \y_{1}^\top\\
                    \x_{1} & {\distmatrix(G)}' & \x_{1}\one^\top + \one \y_{1}^\top\\
                    \y_{1} & \one \x_{1}^\top + \y_{1}\one^\top & {\distmatrix(H)}'
                \end{bmatrix}.
            \end{align}
        since for $u\in V(G)\setminus\{u_{0}\}$ and $v\in V(H)\setminus\{v_{0}\}$ one has
        $d_{G\,\Circ\,H}(u, v)=d_G(u,u_{0})+d_H(v_{0},v)$. Now let $\x\in\mathbb{R}^{1+s+t}$ be the vector given by
            \begin{align}\label{eq:xw}
                \x &= \begin{bmatrix}
                    x_w\\ \y_G \\ \y_H
                \end{bmatrix}, \quad x_w\in\mathbb{R},\, y_G\in\mathbb{R}^{s},\,y_H\in\mathbb{R}^{t}.
            \end{align}
        To guarantee that $\one^\top x = 1$, we put
            \begin{align*}
                x_w = 1-\one^\top \y_G - \one^\top \y_H = 1 - (1-\alpha_G) - (1-\alpha_H) = \alpha_G + \alpha_H - 1.
            \end{align*}
        Then we have that
            \begin{align*}
                \distmatrix \x &= \begin{bmatrix}
                    0 & \x_{1}^\top & \y_{1}^\top\\
                    \x_{1} & {\distmatrix(G)}' & \x_{1}\one^\top + \one \y_{1}^\top\\
                    \y_{1} & \one \x_{1}^\top + \y_{1}\one^\top & {\distmatrix(H)}'
                \end{bmatrix}\begin{bmatrix}
                    x_w\\ \y_G \\ \y_H
                \end{bmatrix}\\
                &= \begin{bmatrix}
                    \x_{1}^\top \y_G + \y_{1}^\top \y_H\\
                    \x_{1} x_w + {\distmatrix(G)}' \y_G + \x_1\one^\top\y_H + \one \y_{1}^\top \y_H\\
                    \y_{1} x_w + \one \x_{1}^\top \y_G + \y_1\one^\top \y_G + {\distmatrix(H)}' \y_H
                \end{bmatrix}\\
                &= \begin{bmatrix}
                    \iota(G) + \iota(H)\\
                    \x_1\,x_w + \bigl(\iota(G)\one - \alpha_G\x_1\bigr) + \x_1(\one^\top\y_H) + \one(\y_1^\top\y_H)\\
                    \y_1\,x_w + \one(\x_1^\top\y_G) + \y_1(\one^\top\y_G) + \bigl(\iota(H)\one - \alpha_H\y_1\bigr)\\
                \end{bmatrix}.
            \end{align*}
        since, by construction, $\alpha_G \x_1 + {\distmatrix(G)}' \y_G = \iota(G)\one$, and similarly $\alpha_H \y_1 + {\distmatrix(H)}' \y_H = \iota(H)\one$. Continuing, we have
            \begin{align*}
                \begin{bmatrix}
                    \iota(G) + \iota(H)\\
                    \x_1\,x_w + \bigl(\iota(G)\one - \alpha_G\x_1\bigr) + \x_1(\one^\top\y_H) + \one(\y_1^\top\y_H)\\
                    \y_1\,x_w + \one(\x_1^\top\y_G) + \y_1(\one^\top\y_G) + \bigl(\iota(H)\one - \alpha_H\y_1\bigr)\\
                \end{bmatrix} &= \begin{bmatrix}
                    \iota(G) + \iota(H)\\
                    \bigl(\iota(G)+\iota(H)\bigr)\one \;+\; \x_1\,(x_w + \one^\top\y_H - \alpha_G)\\
                    \bigl(\iota(G)+\iota(H)\bigr)\one \;+\; \y_1\,(x_w + \one^\top\y_G - \alpha_H)\\
                \end{bmatrix}\\
                &= (\iota(G) + \iota(H))\one,
            \end{align*}
        as claimed, since
            \begin{align*}
                x_w + \one^\top\y_H - \alpha_G &= (\alpha_G + \alpha_H - 1) + (1 - \alpha_H) - \alpha_G = 0,
            \end{align*}
        and similarly for the other term. This proves the first part of the theorem.

        Next, we assume without loss of generality that $\iota(G) = \infty$. In this case, we claim that $\iota(G\,\Circ\,H) = \infty$. By assumption there exists some $\y_G^\top = [\alpha\; \y_G'^\top]$ such that $\distmatrix(G)\y_G = \one$ and $\one^\top \y_G = \alpha + \one^\top \y_G' = 0$.  Define
            \begin{align*}
                \x &= \begin{bmatrix}
                    \alpha \\ \y_G' \\ \vec{0}
                \end{bmatrix} \in\mathbb{R}^{V(G\,\Circ\,H)}.
            \end{align*}
        Then, re-using the block decomposition~\cref{eq:d-decomp}, we have
            \begin{align*}
                \distmatrix \x &= \begin{bmatrix}
                    0 & \x_{1}^\top & \y_{1}^\top\\
                    \x_{1} & {\distmatrix(G)}' & \x_{1}\one^\top + \one \y_{1}^\top\\
                    \y_{1} & \one \x_{1}^\top + \y_{1}\one^\top & {\distmatrix(H)}'
                \end{bmatrix}\begin{bmatrix}
                    \alpha \\ \y_G' \\ \vec{0}
                \end{bmatrix} &= \begin{bmatrix}
                    \x_{1}^\top \y_G'\\
                    \x_{1} \alpha + {\distmatrix(G)}' \y_G'\\
                    \y_{1} \alpha + \one \x_{1}^\top \y_G' + \y_1\one^\top\y_G'
                \end{bmatrix} = \begin{bmatrix}
                    1\\
                    \one\\
                    \one
                \end{bmatrix},
            \end{align*}
        since $\x_1 \alpha + {\distmatrix(G)}' \y_G' = \one$ by construction, and since
            \begin{align*}
                \y_{1} \alpha + \one \x_{1}^\top \y_G' + \y_1\one^\top\y_G' &= \y_{1} \alpha + \one + \y_1(-\alpha) = \one.
            \end{align*}
        But then $\one^\top \x = \alpha + \one^\top \y_G' + 0 = 0$. Thus we have shown that there exists some $\x$ with $\distmatrix(G\,\Circ\,H)\x = \one$ and $\one^\top\x = 0$. Let $\y\in\mathbb{R}^{V(G\,\Circ\,H)}$ be arbitrary such that $\distmatrix(G\,\Circ\,H)\y = \one$. Then
            \begin{align*}
                \distmatrix(G\,\Circ\,H)(\y - \x) = \vec{0},
            \end{align*}
        and thus $\y-\x \in \mathrm{ker}(\distmatrix(G\,\Circ\,H))$. Since $\distmatrix(G\,\Circ\,H)$ is symmetric and $\one\in\mathrm{range}(\distmatrix(G\,\Circ\,H))$, it holds that $\mathrm{ker}(\distmatrix(G\,\Circ\,H)) \perp \one$. Therefore
            \begin{align*}
                \one^\top \y &= \one^\top \x + \one^\top (\y - \x) = 0 + 0 = 0,
            \end{align*}
        from which it follows that $\iota(G\,\Circ\,H) = \infty$.
    \end{proof}

\end{document}
